\documentclass[11pt]{article}
\usepackage[utf8]{inputenc}
\usepackage[T1]{fontenc}
\usepackage{amsmath,amssymb}
\usepackage{graphicx}
\usepackage{booktabs}
\usepackage[margin=1in]{geometry}
\usepackage{hyperref}

\title{Arboles de Decision: Construyendo Intuicion\\
\large (Notebook: large\_trees\_exploration)}
\author{}
\date{}

\begin{document}
\maketitle

\begin{abstract}
Este documento explica, paso a paso, como funcionan las estrategias de
pooled testing usando arboles de decision.
Empezamos con 2 personas y 1 test, y vamos subiendo hasta 8 personas
con poblaciones heterogeneas.
No se cubren aqui los parametros alpha/beta (ver documentos 02 y 03)
ni las comparaciones con la tesis de Nico (ver documento 03).
\end{abstract}

%% ====================================================================
\section{Caso mas simple: 2 personas, 1 test}

Tenemos 2 personas, cada una con probabilidad $p=0.2$ de estar infectada.
Utilidad $u=1$ por persona aclarada. Solo 1 test disponible.

\textbf{Opcion A:} testear persona 0 sola.
\begin{itemize}
  \item Si sale negativo (prob 0.8): aclaramos 1 persona. EU = 0.8.
  \item Si sale positivo: no aclaramos a nadie.
\end{itemize}

\textbf{Opcion B:} testear a las 2 juntas (pool).
\begin{itemize}
  \item Si $r=0$ (prob $0.8 \times 0.8 = 0.64$): aclaramos a las 2. EU = $0.64 \times 2 = 1.28$.
  \item Si $r>0$: no aclaramos a nadie.
\end{itemize}

\textbf{Resultado:} agrupar gana (EU 1.28 vs 0.80) porque con baja prevalencia
es probable que ambas esten sanas.

%% ====================================================================
\section{Efecto del valor: utilidades asimetricas}

Misma configuracion pero ahora persona A vale $u_A = 10$, persona B vale $u_B = 1$.
Ambas con $p = 0.1$.

Cuando una persona vale mucho mas, puede convenir testearla sola
para no ``desperdiciar'' el test si la otra esta infectada.

Al barrer $u_A$ de 1 a 20, hay un punto de cruce donde individual
empieza a ganarle al pool.

%% ====================================================================
\section{Efecto de la prevalencia}

$n=3$, $B=2$, $G=3$. Tres personas iguales ($u=1$), variamos $p$:

\begin{itemize}
  \item $p=0.05$ (baja): conviene pool grande. Casi seguro todos sanos.
  \item $p=0.20$ (media): pool sigue ganando, pero con menos margen.
  \item $p=0.50$ (alta): pools chicos o individuales. Muchos positivos.
\end{itemize}

\textbf{Regla general:} baja prevalencia $\rightarrow$ pools grandes.
Alta prevalencia $\rightarrow$ pools chicos.

%% ====================================================================
\section{Aumentado vs Clasico}

$n=3$, $B=2$, $G=3$, probabilidades $p = [0.2, 0.25, 0.15]$, utilidades $u = [3, 2, 4]$.

\begin{itemize}
  \item \textbf{Test clasico:} solo dice positivo/negativo (binario).
  \item \textbf{Test aumentado:} dice \emph{cuantos} infectados hay ($r = 0, 1, 2, \ldots$).
\end{itemize}

El aumentado da estrictamente mas informacion.
Permite mejores actualizaciones bayesianas despues de un resultado positivo.
Por ejemplo, si $r=1$ en un pool de 3, sabemos que exactamente 1 esta infectado.
El clasico solo sabe que ``al menos 1'' lo esta.

\textbf{Resultado:} aumentado siempre $\geq$ clasico en utilidad esperada.

%% ====================================================================
\section{Poder del presupuesto}

$n=4$, $G=4$, $p=0.15$, $u=1$. Variamos $B = 1, 2, 3, 4$.

Mas tests $=$ mas oportunidades de aclarar personas.
La utilidad esperada sube con $B$ pero con rendimientos decrecientes:
cada test adicional aporta menos que el anterior.
Con suficientes tests, nos acercamos a $U_{\max}$ (la cota superior si
supieramos quien esta infectado).

%% ====================================================================
\section{Por que no testear uno por uno?}

$n=4$, $B=2$, $p=0.15$, $u=1$. Variamos $G = 1, 2, 3, 4$.

$G=1$ obliga a tests individuales: con $B=2$ solo aclaramos 2 personas.
$G=4$ permite meter a las 4 en un pool: si sale $r=0$, aclaramos a las 4 de golpe.

\textbf{Resultado:} a mayor $G$, mejor EU (cuando la prevalencia es baja).
Agrupar multiplica el impacto de cada test.

%% ====================================================================
\section{Poblacion heterogenea: VIPs vs comunes}

$n=8$, $B=3$, $G=3$.
\begin{itemize}
  \item 4 VIPs: $p=0.3$, $u=10$ (alto valor, moderada infeccion).
  \item 4 comunes: $p=0.1$, $u=1$ (bajo valor, baja infeccion).
\end{itemize}

El optimizador (DP) prioriza aclarar a los VIPs porque valen 10x mas.
El greedy puede no hacerlo bien porque es miope (solo ve un paso adelante).

Al variar $p_{\text{VIP}}$ de 0.1 a 0.5, el gap greedy--optimo crece:
a mayor heterogeneidad, mas importa optimizar bien.

%% ====================================================================
\section{Busqueda binaria como heuristica}

$n=6$, $B=4$, $G=3$, $p=0.15$, $u=10$.

La busqueda binaria es una heuristica natural: testea la mitad, si sale positivo
subdivide. Pero el DP optimo es mas flexible --- puede elegir pools
que no son mitades exactas, explotando la informacion del conteo $r$.

Comparacion:
\begin{itemize}
  \item DP optimo: mejor EU.
  \item Greedy: casi tan bueno como DP.
  \item Busqueda binaria: peor que ambos (no usa info de conteo).
\end{itemize}

%% ====================================================================
\section{Greedy vs Optimo: cuanto se pierde?}

Grid de instancias aleatorias: $n \in \{4,5,6\}$, $B \in \{2,3\}$, $G \in \{3,4\}$.

El gap greedy--optimo varia por instancia (0--15\%).
\textbf{Greedy counting} (usa toda la historia bayesiana) supera a
\textbf{greedy secuencial} (actualiza test por test).

%% ====================================================================
\section{Solver hibrido: K pasos greedy + DP}

$n=6$, $B=4$, $G=3$. Variamos $K$ de 0 (todo DP) a $B$ (todo greedy).

\begin{itemize}
  \item $K=0$: solucion exacta, pero lenta.
  \item $K=B$: rapido, pero pierde calidad.
  \item $K=2$--$3$: recupera 80--90\% de la calidad del DP con una fraccion del tiempo.
\end{itemize}

Este es el enfoque practico: usar greedy para las primeras decisiones
(donde el espacio es grande) y DP para las ultimas (donde el espacio
ya se redujo).

\end{document}
